\section{Singularly foliated bundles and their holonomy groupoids}\label{sec4}

\subsection{Singularly foliated bundles}

Singular foliations are generalisations of regular foliations, and are naturally associated to Lie groupoids more generally. In~each of these special cases, one has a notion of fibre bundle which is compatible with the additional structure~-- in the case of a Lie groupoid action, the correct notion is that of an \textit{equivariant bundle}, while for a regular foliation the correct notion is that of a \textit{foliated bundle} in the sense of Kamber and Tondeur~\cite{folbund}. We~give in this section what appears to be the first definition of a fibre bundle compatible with a singular foliation, which simultaneously generalises equivariant and foliated bundles.

First, notice that projectable vector fields on a fibre bundle $\pi_{B}\colon B\rightarrow M$ over a manifold $M$ \textit{do not} generally form a sheaf of $C^{\infty}_{B}$-modules over $B$. Indeed, if $X$ is any projectable vector field and $f\in C^{\infty}(B)$ is any function which is non-constant along the fibres of $\pi_{B}$, then equation~\eqref{projectableeqn} will in general no longer hold for the vector field $fX$. Projectable vector fields are, however, closed under multiplication by functions of the form $f\circ\pi_{B}$, where $f\in C^{\infty}(M)$. Since $\pi_{B}$ is an open map, we can formulate the following definition, which will play a crucial role in our definition of singularly foliated bundle.

\begin{Definition}%\label{projectsheaf}
Let $\pi_{B}\colon B\rightarrow M$ be a fibre bundle. Denote by $C^{\infty}_{\proj,B}$ the subsheaf
\begin{gather*}
C^{\infty}_{\proj,B}(\OO):=\big\{f\circ\pi_{B}\in C^{\infty}_{B}(\OO)\colon f\in C^{\infty}_{M}(\pi_{B}(\OO))\big\}
\end{gather*}
of $C^{\infty}_{B}$, which we call the \textit{sheaf of projectable functions}. We~denote by $\XF_{\proj,B}$ the sheaf of $C^{\infty}_{\proj,B}$-modules
\begin{gather*}
\XF_{\proj,B}(\OO)\!:=\!\big\{X{\in}\XF_{B}(\OO)\colon \text{there is $(\pi_{B})_{*}X{\in}\XF_{M}(\pi_{B}(\OO))$ with ${\rm d}\pi_{B}\!\circ X\! = \! (\pi_{B})_{*}(X)\!\circ\!\pi_{B}$}\big\},
\end{gather*}
which we call the \textit{sheaf of projectable vector fields}.
\end{Definition}

The pushforward of projectable vector fields can now be characterised in the following sheaf-theoretic fashion. Recall that for a map $f\colon X\rightarrow Y$ of topological spaces and a sheaf $\SS$ on $X$, we use $f_{!}\SS$ to denote the pushforward of $\SS$ on $Y$~\cite[p.~65]{hartshorne}.

\begin{Proposition}
Let $\pi_{B}\colon B\rightarrow M$ be a fibre bundle. The~pushforward of projectable vector fields induces a morphism $(\pi_{B})_{*}\colon (\pi_{B})_{!}\XF_{\proj,B}\rightarrow\XF_{M}$ of sheaves of $C^{\infty}_{M}$-modules that preserves the Lie bracket.
\end{Proposition}

\begin{proof}
Notice first that we have a canonical isomorphism $C^{\infty}_{M}\cong (\pi_{B})_{!}C^{\infty}_{\proj,B}$ of sheaves of rings, obtained simply by sending $f\in C^{\infty}_{M}(\OO)$ to $f\circ\pi_{B}\in C^{\infty}_{\proj,B}(\pi_{B}^{-1}(\OO))$ for each open set $\OO$ in $M$. In~this way the $C^{\infty}_{\proj,B}$-module structure of $\XF_{\proj,B}$ indeed defines a $C^{\infty}_{M}$-module structure on $(\pi_{B})_{!}\XF_{\proj,B}$. The~pushforward $(\pi_{B})_{*}$ of $X\in\XF_{\proj,B}(\pi_{B}^{-1}(\OO))$ to $(\pi_{B})_{*}(X)\in\XF_{M}(\OO)$ then clearly preserves the associated $C^{\infty}_{M}(\OO)$-module structure for each open set $\OO$, and it is well-known~\cite[Lemma~3.10]{natopdiffgeom} that it also preserves the Lie bracket of vector fields.
\end{proof}

Singularly foliated bundles are now defined by \textit{singular partial connections}, which are particularly well-behaved partially-defined right-inverses of the pushforward morphism.

\begin{Definition}\label{singfolbund}
A \textit{singularly foliated bundle} is a triple $(\pi_{B},\FF,\ell)$, where $\pi_{B}\colon B\rightarrow M$ is a fibre bundle, $\FF$ is a singular foliation of $M$, and $\ell\colon \FF\rightarrow(\pi_{B})_{!}\XF_{\proj,B}$ is a morphism of sheaves of $C^{\infty}_{M}$-modules which preserves the Lie bracket, and for which the following hold.
\begin{enumerate}\itemsep=0pt
\item The morphism $\ell$ is a \textit{partial right-inverse} to $(\pi_{B})_{*}$ in the sense that
\begin{gather*}
(\pi_{B})_{*}\circ\ell = \id_{\FF},
\end{gather*}
on the sheaf $\FF$. In~particular this implies that $\ell$ is injective.
\item The morphism $\ell$ is \textit{complete} in the sense that for any open set $\OO$ in $M$, any $X\in\FF(\OO)$ and any $x\in\OO$, if $\Fl^{X}(x)$ is defined on an interval $I\subset\RB$ then so too is $\Fl^{\ell(X)}(b)$ for any $b\in B_{x}$.
\item The morphism $\ell$ is \textit{smooth} in the sense that the induced morphism
\begin{gather*}
\Gamma_{\loc}(\FF)\rightarrow \Gamma_{\loc}\big((\pi_{B})_{!}\XF_{\proj,B}\big)
\end{gather*}
of diffeological spaces is smooth with respect to the diffeology of Proposition~\ref{functprop}.
\end{enumerate}
We refer to such a morphism $\ell$ as a \textit{singular partial connection}.
\end{Definition}

We will usually denote a singularly foliated bundle $(\pi_{B},\FF,\ell)$ by simply $\pi_{B}$, with $\FF$ and $\ell$ assumed unless otherwise stated. Before discussing some examples, let us mention that completeness of $\ell$ does \textit{not} automatically follow from $\ell$ being a partial right inverse to $(\pi_{B})_{*}$. Indeed, it is easy to verify that for any open set $\OO\subset M$ and for any $X\in\FF(\OO)$, we have the relationship
\begin{gather*}
\Fl^{X}_{t}(x) = \pi_{B}\big(\Fl^{\ell(X)}_{t}(b)\big),\qquad x\in\OO,\ b\in B_{x}
\end{gather*}
between the flows of $\Fl^{X}(x)$ and $\Fl^{\ell(X)}(b)$ wherever they are defined. In~particular, that $\ell$ is a partial right-inverse to $(\pi_{B})_{*}$ implies that for any $b\in B_{x}$, the domain of $\Fl^{\ell(X)}(b)$ is contained in the domain of $\Fl^{X}(x)$. The~converse, however, does \textit{not} follow without completeness of $\ell$, as is easily seen by considering the standard example of $B =\RB^{2}$, $M = \RB$, and with $\ell\colon \XF(M)\rightarrow\XF_{\proj,B}(B)$ defined by
\begin{gather*}
\ell\bigg(f\frac{\partial}{\partial x}\bigg)(x,y):=f(x)\bigg(\frac{\partial}{\partial x}+y^{2}\frac{\partial}{\partial y}\bigg)
\end{gather*}
for $f\in C^{\infty}(M)$ and $(x,y)\in B$.

\begin{Example}[trivial bundles]\label{triv}
If $(M,\FF)$ is any singularly foliated manifold and $Q$ is any other manifold, then the trivial bundle $\pi\colon M\times Q\rightarrow M$ is canonically a singularly foliated bundle. Indeed, with respect to the decomposition $T(M\times Q)\cong TM\times TQ$, one has the trivial lift $\ell\colon \XF_{M}\rightarrow\pi_{!}\XF_{\proj,M\times Q}$ defined by the formula
\begin{gather*}
\ell(X):=(\pi^{*}(X),0)\in\XF_{M\times Q}\big(\pi^{-1}(\OO)\big) = (\pi_{!}\XF_{M\times Q})(\OO)
\end{gather*}
for all open sets $\OO$ in $M$ and $X\in\XF_{M}(\OO)$. Restricting $\ell$ to the subsheaf $\FF$ of $\XF_{M}$ one obtains a singular partial connection, with both completeness and smoothness being trivial.
\end{Example}

\begin{Example}[regularly foliated bundles]\label{regex}
Suppose that $\pi_{B}\colon B\rightarrow M$ is a \textit{regularly} foliated bundle, in the sense of Kamber--Tondeur~\cite[Definition~2.1]{folbund}~-- that is, there exists an involutive subbundle $T\FF_{B}\subset TB$ which is projected fibrewise-injectively to a subbundle $T\FF_{M}$ of $TM$. Involutivity of $T\FF_{B}$ implies that both $T\FF_{B}$ and $T\FF_{M}$ integrate to regular foliations of $B$ and $M$ respectively. Now if $\OO$ is any open subset of $M$ and $X\in\FF_{M}(\OO)$, then one obtains $\ell(X)\in\XF_{\proj,B}\big(\pi_{B}^{-1}(\OO)\big)$ whose value at a point $b\in B$ is the unique vector in $T_{b}\FF_{B}$ that is mapped by ${\rm d}\pi_{B}$ to $X(\pi_{B}(b))$. Clearly then the resulting morphism $\ell\colon \FF_{M}\rightarrow(\pi_{B})_{!}\XF_{\proj,B}$ is a singular partial connection in the sense of Definition~\ref{singfolbund}. Completeness and smoothness can both be seen by choosing foliated coordinates, in which $\ell$ is simply given by a trivial lift as in~Example~\ref{triv}.

Conversely, suppose that $\FF$ is a regular foliation of a manifold $M$, with leaf dimension $p$, and that $\pi_{B}\colon B\rightarrow M$ is a fibre bundle with a singular partial connection $\ell\colon \FF\rightarrow(\pi_{B})_{!}\XF_{\proj,B}$. In~a foliated chart $\OO\cong\RB^{p}\times\RB^{q}$ of $M$, wherein $\FF(\OO)$ is the $C^{\infty}_{M}(\OO)$-span of vector fields $\{e_{1},\dots,e_{p}\}$ that form the standard frame field of $\RB^{p}$, injectivity of $\ell$ implies that the vector fields $\{\ell(e_{1}),\dots,\ell(e_{p})\}$ span a $p$-dimensional subspace of $T_{b}B$ at each point $b\in\pi_{B}^{-1}(\OO)$ which intersects the vertical tangent space at $b$ only through zero. One thus obtains a smooth $p$-dimensional distribution $T\FF_{B}$ in $B$, and involutivity of $\FF$ together with the fact that $\ell$ preserves the Lie bracket implies that $T\FF_{B}$ is involutive. Thus $\pi_{B}\colon B\rightarrow M$ is a foliated bundle in~the sense of Kamber--Tondeur.
\end{Example}

\begin{Example}[equivariant bundles]\label{equiex}
Let $\GG$ be a Lie groupoid with unit space $M$. Denote the Lie algebroid of $\GG$ by $A:=\ker({\rm d}s)|_{M}$, with anchor map ${\rm d}r\colon A\rightarrow TM$, and let $\FF$ denote the associated sheaf of vector fields of the form ${\rm d}r\circ\sigma$, where $\sigma$ is an element of the sheaf of sections~$\AA$ of the Lie algebroid $A$. Let $\pi_{B}\colon B\rightarrow M$ be a fibre bundle.

If $\GG$ acts on $B$, then denote by $B\rtimes\GG$ the associated action groupoid~\cite[Example 2.2]{aac1}, with range and source denoted $r_{B}$ and $s_{B}$ respectively. Then the Lie algebroid $A_{B}:=\ker({\rm d}s_{B})|_{B}$ associated with the action groupoid $B\rtimes\GG$ is isomorphic to $\pi_{B}^{*}(A)$. For any open subset $\OO$ of~$M$, the formula
\begin{gather*}
\tilde{\ell}({\rm d}r\circ\sigma):={\rm d}r_{B}\circ\pi_{B}^{*}(\sigma),\qquad \sigma\in\AA(\OO)
\end{gather*}
then defines a singular partial connection. Completeness and smoothness are consequences of~the fact that $\tilde{\ell}$ is defined in terms of a smooth action of the groupoid $\GG$.
\end{Example}


\subsection{The holonomy groupoids of singularly foliated bundles}

One of the key objects in our construction of the holonomy groupoids of a singularly foliated bundle are pseudo-bundles of \textit{invariant} germs/jets. To~define these pseudo-bundles, we need \textit{vertical} prolongations of projectable vector fields (cf.~\cite[Definition~1.15]{varbi}).

\begin{Definition}
Let $\pi_{B}\colon B\rightarrow M$ be a fibre bundle, let $X\in\XF_{\proj}(B)$, and let $k$ denote any of~the symbols $1,\dots,\infty,\g$. The~vector field $\vpf^{k}(X)$ on $\Gamma_{k}(\pi_{B})$ defined by
\begin{gather*}
\vpf^{k}(X)\big(x,[\sigma]^{k}_{x}\big):=\frac{\rm d}{{\rm d}t}\bigg|_{t=0}\Big(x,\big[\Fl^{X}_{t}\circ\sigma\circ\Fl^{(\pi_{B})_{*}(X)}_{-t}\big]_{x}^{k}\Big)
\end{gather*}
is called the \textit{vertical $k$-prolongation} of $X$. On $\Gamma_{0}(\pi_{B}) = B$, we set $\vpf^{0}(X):=0$.
\end{Definition}

For a projectable vector field $X$ on a fibre bundle $\pi_{B}\colon B\rightarrow M$ and $k\geq1$, the vertical $k$-prolongation $\vpf^{k}(X)$ is the image of $\pf^{k}(X)$ under the canonical $\big(\pi_{B}^{k-1,k}\big)^{*}\big(V\pi_{B}^{k-1}\big)$-valued contact form $\theta^{(k)}$~\cite[Chapter~6.3]{saunders} on $\Gamma_{k}(\pi_{B})$. In~local coordinates $(x^{i},f^{\alpha},f^{\alpha}_{i},\dots)$, the components of~$\theta^{(k)}$ are given by
\begin{gather*}
\big(\theta^{(k)}\big)^{\alpha}_{I} = {\rm d}f^{\alpha}_{I} - f^{\alpha}_{Ii}{\rm d}x^{i}
\end{gather*}
for each $|I|<k$. Thus it is easily checked (cf.~equation~\eqref{prolongation}) that in coordinates the vertical prolongation of $X = a^{i}\frac{\partial}{\partial x^{i}}+b^{\alpha}\frac{\partial}{\partial f^{\alpha}}$ is given by
\begin{gather*}
\vpf^{k}(X) = \sum_{|I|=0}^{k-1}D_{I}^{(k)}(b^{\alpha}-a^{i}f^{\alpha}_{i})\frac{\partial}{\partial f^{\alpha}_{I}}.
\end{gather*}
The invariant pseudo-bundles of a singularly foliated bundle are now defined as follows.

\begin{Definition}\label{FFinv}
Let $\pi_{B}\colon B\rightarrow M$ be a singularly foliated bundle, and let $k$ denote any of the symbols $0,1,\dots,\infty,\g$. An element $\big(x,[\sigma]^{k}_{x}\big)$ of $\pi_{B}$ is said to be \textit{$\FF$-invariant to order $k$} if
\begin{gather*}
\vpf^{k}\big(\ell(X)\big)\big(x,[\sigma]^{k}_{x}\big) = 0
\end{gather*}
for all $X\in\FF$ defined in a neighbourhood $x$. We~say that $\pi_{B}$ \textit{admits enough conservation laws to order $k$} if at each $x\in M$ there is $\big(x,[\sigma]^{k}_{x}\big)$ which is $\FF$-invariant. This being the case, we call the corresponding sub-pseudo-bundle $\Gamma_{k}(\pi_{B})^{\FF}\subset\Gamma_{k}(\pi_{B})$ consisting of invariant germs/jets the \textit{$k^{\rm th}$ order $\FF$-invariant pseudo-bundle} of $\pi_{B}$. We~denote by $\pi_{B}^{k,\FF}$ the restriction of $\pi_{B}^{k}$ to $\Gamma_{k}(\pi_{B})^{\FF}$.
\end{Definition}

Note that if a singularly foliated bundle admits enough conservation laws to order $k>0$, then it admits enough conservation laws to any $l\leq k$. Moreover, by definition, \textit{every} singularly foliated bundle admits enough conservation laws to order 0.

Recall now~\cite[Definition~2.10]{mac3} that if $\pi_{B}\colon B\rightarrow M$ is a \textit{regularly} foliated bundle, a locally-defined section $\sigma$ of~$\pi_{B}$ is said to be \textit{distinguished} if, about any point in its image, there exist foliated coordinates $(x_{\alpha},y_{\alpha},f_{\alpha})$, with $x_{\alpha}$ and $y_{\alpha}$ denoting the leafwise and transverse coordinates respectively in the base, and with $f_{\alpha}$ denoting coordinates in the fibre, with respect to which $\sigma = \sigma(y_{\alpha})$ is independent of~the leafwise coordinates. We~denote by $\DS_{\g}(\pi_{B})$ the diffeological bundle of~germs of~the sheaf of~distinguished sections, and by $\DS_{k}(\pi_{B})$ the bundle of~jets of~distinguished sections. The~next proposition says that the $\FF$-invariant pseudo-bundles of~Defi\-ni\-tion~\ref{FFinv} generalise the bundles of~distinguished sections appearing in the regular case, and that therefore our constructions recover those of~\cite{mac3} in the regular case. In~particular, since distinguished functions of~this sort furnish the (local) degree 0 characteristic cohomology classes of~regular foliations~\cite[Example 1]{bryantgriffiths1}, we feel justified in referring to the $\FF$-invariant germs/jets of~a singularly foliated bundle as its \textit{conservation laws}.

\begin{Proposition}%\label{reg}
Let $\pi_{B}\colon B\rightarrow M$ be a regularly foliated bundle, and let $k$ denote any of the symbols $0,\dots,\infty,\g$. Then the diffeological subspace $\Gamma_{k}(\pi_{B})^{\FF}$ coincides with the space $\DS_{k}(\pi_{B})$ of classes of distinguished sections.
\end{Proposition}

\begin{proof}
In foliated coordinates $(x_{\alpha},y_{\alpha},f_{\alpha})$ for $B$, corresponding to leafwise, transverse and fibre coordinates respectively, any element $X\in\FF$ is given by some $C^{\infty}_{M}$-linear combination
\begin{gather*}
X = a^{i}\frac{\partial}{\partial x_{\alpha}^{i}},
\end{gather*}
while $\ell(X)$ (see Example~\ref{regex}) is given by
\begin{gather*}
\ell(X) = \pi_{B}^{*}(a^{i})\frac{\partial}{\partial x_{\alpha}^{i}}.
\end{gather*}
Thus, in our coordinates $(x_{\alpha},y_{\alpha},f_{\alpha})\in\RB^{p}\times\RB^{q}\times\RB^{k}$, we have simply
\begin{gather*}
\Fl^{\ell(X)}_{t}(x_{\alpha},y_{\alpha},f_{\alpha}) = \big(\Fl^{X}_{t}(x_{\alpha}),y_{\alpha},f_{\alpha}\big)
\end{gather*}
for small $t$. It follows immediately that for any smooth function $\sigma\colon \RB^{p}\times\RB^{q}\rightarrow\RB^{k}$, the curve
\begin{gather*}
t\mapsto\big(\Fl^{\ell(X)}_{t}\circ(\id\times\sigma)\circ\Fl^{X}_{-t}\big)(x_{\alpha},y_{\alpha}) = \big(x_{\alpha},y_{\alpha},\sigma\big(\Fl^{X}_{t}(x_{\alpha}),y_{\alpha}\big)\big)
\end{gather*}
is constant in $t$ \textit{for all $X\in\FF$} if and only if $\sigma$ is constant in the $x$ coordinate. Thus $\Gamma_{\g}(\pi_{B})^{\FF} = \DS_{\g}(\pi_{B})$. Since foliated coordinates can always be used to extend an invariant jet to an invariant local section, one has $\Gamma_{k}(\pi_{B})^{\FF} = \DS_{k}(\pi_{B})$ for $k\leq\infty$ also.
\end{proof}

Defining lifting maps and leafwise transport functors for a singularly foliated bundle is now a simple matter of putting our definitions together.

\begin{Theorem}\label{liftsmooth}
Let $\pi_{B}\colon B\rightarrow M$ be a singularly foliated bundle. Let $k$ denote any of the symbols $0,\dots,\infty,\g$, and suppose that $\pi_{B}$ admits enough conservation laws to order $k$. Then for each $\big(\gamma,[X]^{\g},d\big)\in\PP(\FF)$, $\gamma(0) = x$, and for each $\big(x,[\sigma]^{k}_{x}\big)\in\Gamma_{k}(\pi_{B})^{\FF}_{x}$, the map
\begin{gather}\label{ivpsol}
t\mapsto\big(\gamma(t),\big[\Fl^{\ell(X)}_{t,0}\circ\sigma\circ\Fl^{X}_{0,t}\big]^{k}_{\gamma(t)}\big)
\end{gather}
is the unique solution to the initial value problem
\begin{gather}\label{ivp}
\frac{\rm d}{{\rm d}s}\bigg|_{t}f\big(s;x,[\sigma]^{k}_{x}\big) = \pf^{k}\big(\ell(X(t))\big)\big(f\big(t;x,[\sigma]^{k}_{x}\big)\big),\qquad
f\big(0;x,[\sigma]^{k}_{x}\big) = \big(x,[\sigma]^{k}_{x}\big)
\end{gather}
in the diffeological space $\Gamma_{k}(\pi_{B})^{\FF}$ for which $\pi_{B}^{k,\FF}\circ f\big({-};x,[\sigma]^{k}_{x}\big) = \gamma(-)$. Moreover, the lifting map $L\big(\pi_{B}^{k,\FF}\big)\colon \PP(\FF)\times_{s,\pi_{B}^{k,\FF}}\Gamma_{k}(\pi_{B})^{\FF}\rightarrow \PP\big(\Gamma_{k}(\pi_{B})^{\FF}\big)$ defined by
\begin{gather}\label{Leqn}
L\big(\pi_{B}^{k,\FF}\big)\big(\gamma,[X]^{\g},d;x,[\sigma]^{k}_{x}\big)(t):=
\big(\gamma(t),\big[\Fl^{\ell(X)}_{t,0}\circ\sigma\circ\Fl^{X}_{0,t}\big]^{k}_{\gamma(t)}\big),\qquad t\in[0,\infty)
\end{gather}
is smooth.
\end{Theorem}

Note that the expression on the right hand side of equation~\eqref{Leqn} only makes sense by the completeness assumption on the singular partial connection. To~prove Theorem~\ref{liftsmooth}, we require the following lemma.

\begin{Lemma}\label{technicallemma}
Let $(M,\FF)$ be a singular foliation. Suppose that $U$ is an open set in Euclidean space and $X\colon \RB_{\geq0}\times U\rightarrow\Gamma_{\loc}(\FF)$ and $\gamma\colon \RB_{\geq0}\times U\rightarrow M$ are smooth maps for which
\begin{enumerate}\itemsep=0pt
\item[$(1)$] $\dom(X(t,u)) = \dom(X(s,u))$ for all $s,t\in\RB_{\geq0}$ and $u\in U$, and for which
\item[$(2)$] $t\mapsto\gamma(t,u)$ is an integral curve of $X(t,u)$ for all $(t,u)\in\RB_{\geq0}\times U$. That is,
\begin{gather*}
\frac{\rm d}{{\rm d}s}\bigg|_{s=t}\gamma(s,u) = X(t,u)(\gamma(t,u))
\end{gather*}
for all $(t,u)\in\RB_{\geq0}\times U$.
\end{enumerate}
Then for each $t_{0}\in\RB_{\geq0}$ and $u_{0}\in U$, there exist $\epsilon>0$ and open neighbourhoods $\UU\ni u_{0}$ in $U$ and $\OO\ni x_{0}:=\gamma(0,u_{0})$ in $M$, such that $\OO\subset\dom\big(\Fl^{X(u)}_{t,0}\big)$ for all $(t,u)\in[0,t_{0}+\epsilon)\times\UU$, and for which the map
\begin{gather*}
[0,t_{0}+\epsilon)\times\UU\times\OO\ni(t,u,x)\mapsto\Fl^{X(u)}_{t,0}(x)\in M
\end{gather*}
is smooth.
\end{Lemma}

\begin{proof}
By the definition of the diffeology on $\Gamma_{\loc}(\FF)$ and hypothesis 1, we can find open sets $\UU'\ni u_{0}$ and $\OO'\ni x_{0}$ such that $\OO'\subset\dom(X(t,u))$ for all $(t,u)\in\RB_{\geq0}\times\UU'$, and on which
\begin{gather*}
\RB_{\geq0}\times \UU'\times\OO'\ni(t,u,x)\mapsto X(t,u)(x)\in TM
\end{gather*}
is smooth. By hypothesis 2 we can assume that $\OO'$ contains $\gamma(\RB_{\geq0}\times\{u_{0}\})$. Now $X$ may be regarded as a time-dependent vector field on the manifold $\UU'\times\OO'$, and then by standard theory~\cite[Theorem~9.48]{leesmth}, there exists some maximal open neighbourhood $\NN\subset\UU'\times\OO'$ of~$(u_{0},x_{0})$ on which $t\mapsto\Fl^{X(u)}_{t,0}(x)$ is defined for all $u\in\UU$, $x\in\OO$ and for small $t$. Since by hypothesis $t\mapsto\Fl^{X(u_{0})}_{t,0}(x_{0}) = \gamma(t,u_{0})$ is defined for all $t\in\RB_{\geq0}$ we can always choose $\epsilon>0$, $\UU\ni u_{0}$ and $\OO\ni x_{0}$ small enough that $(t,u,x)\mapsto\Fl^{X(u)}_{t,0}(x)$ is well-defined and smooth on $[0,t_{0}+\epsilon)\times\UU\times\OO$ as claimed.
\end{proof}

\begin{proof}[Proof of Theorem~\ref{liftsmooth}]
That equation~\eqref{ivpsol} defines a curve in $\Gamma_{k}(\pi_{B})^{\FF}$ follows from the inva\-riance of $\FF$ under its own (possibly time-dependent~\cite[Lemma~3.3]{VilGar1}) flows~\cite[Proposition~1.6]{iakovos2}, the $\FF$-invariance of $\sigma$, and that $\ell$ is bracket preserving. More precisely, given $t\in[0,\infty)$, an open neighbourhood $\OO$ of $\gamma(t)$ and $Y\in\FF(\OO)$, since $\ell$ is bracket preserving we have
\begin{gather}\label{bracketpres}
\ell([X(t),Y])\big(\Fl^{\ell(X)}_{t,0}(b)\big) = [\ell(X(t)),\ell(Y)]\big(\Fl^{\ell(X)}_{t,0}(b)\big)
\end{gather}
for all $b\in\Fl^{\ell(X)}_{0,t}(\OO)$. Now, for any such $b$, with $\pi_{B}(b) = x'$, the right hand side of equation~\eqref{bracketpres} is the lift by $\ell$ of the tangent to the curve $r\mapsto\big(\Fl^{X}_{0,t+r}\big)_{*}\big(Y\big(\Fl_{t+r,0}^{X}(x')\big)\big)$ at $r = 0$,
while the left hand side is the tangent to the curve $r\mapsto \big(\Fl^{\ell(X)}_{0,t+r}\big)_{*}\big(\ell(Y)\big(\Fl^{\ell(X)}_{t+r,0}(b)\big)\big)$ at $r=0$. By uniqueness of flows therefore, we have
\begin{gather}\label{ellhom}
\big(\Fl^{\ell(X)}_{0,t}\big)_{*}(\ell(Y)) = \ell\big(\big(\Fl^{X}_{0,t}\big)_{*}(Y)\big).
\end{gather}
For notational simplicity denote $\varphi:=\Fl^{X}_{t,0}$ and $\ell(\varphi):=\Fl^{\ell(X)}_{t,0}$. Then we use equation~\eqref{ellhom} to compute
\begin{align*}
\vpf^{k}(\ell(Y))\big(\gamma(t),[\ell(\varphi)\!\circ\!\sigma\!\circ\!\varphi^{-1}]^{k}_{\gamma(t)}\big)
&= \frac{\rm d}{{\rm d}s}\bigg|_{s=0}\!\!\big(\gamma(t),\big[\Fl^{\ell(Y)}_{s}\circ\ell(\varphi) \circ\sigma\circ\varphi^{-1}\circ\Fl^{Y}_{-s}\big]^{k}_{\gamma(t)}\big)
\\
&= \frac{\rm d}{{\rm d}s}\bigg|_{s=0}\!\!\big(\gamma(t),\!\big[\ell(\varphi) \!\circ\Fl^{\ell(\varphi^{-1}_{*}(Y))}_{s}\! \circ\sigma\circ\Fl^{\varphi^{-1}_{*}(Y)}_{-s} \circ\varphi^{-1}\big]^{k}_{\gamma(t)}\big)
\\
&= 0
\end{align*}
by the $\FF$-invariance of $\sigma$ and since $\varphi^{-1}_{*}(Y)\in\FF$. Therefore equation~\eqref{ivpsol} does indeed define a curve in $\Gamma_{k}(\pi_{B})^{\FF}$. That equation~\eqref{ivpsol} defines a solution to the initial value problem of~equation~\eqref{ivp} follows by definition of the prolongation.

Uniqueness for $k$ finite follows from uniqueness of the flow in a manifold (since each $\Gamma_{k}(\pi_{B})^{\FF}$ is a subset of the jet manifold $\Gamma_{k}(\pi_{B})$), while for $k=\infty$ any solution to equation~\eqref{ivp} is the projective limit of the solutions for finite $k$, so uniqueness in this case follows from uniqueness for each finite $k$. Since $\Gamma_{\g}(\pi_{B})$ is neither a manifold nor a projective limit of manifolds, the argument is more subtle in the $k=\g$ case. Suppose that $\rho\colon [0,d]\rightarrow\Gamma_{\g}(\pi_{B})^{\FF}$ is some other solution to the initial value problem given in equation~\eqref{ivp}; thus in particular $\rho(0) = \big(x,[\sigma]^{\g}_{x}\big)$. Let us consider $\epsilon>0$ sufficiently small that for all $t\in[0,\epsilon)$ we can write $\rho(t) = \big(\gamma(t),[\sigma_{t}]^{\g}_{\gamma(t)}\big)$ for some $\FF$-invariant family $t\mapsto\sigma_{t}$ in $\Gamma_{\loc}(\pi_{B})$. By definition of the diffeology on $\Gamma_{\loc}(\pi_{B})$ we may assume that $\epsilon$ is sufficiently small that there is some open neighbourhood $\OO$ of $\gamma(0)$ containing $\gamma([0,\epsilon))$ and on which $\sigma_{t}$ is defined for all $t\in[0,\epsilon)$. Let us also assume that there are fibre bundle coordinates $\big(x^{1},\dots,x^{n};b^{1},\dots,b^{m}\big)$ on $B$ about $\sigma(x)$, the projection of whose domain to~$M$ contains $\OO$. By hypothesis we have\vspace{-1ex}
\begin{gather*}
\frac{\rm d}{{\rm d}s}\bigg|_{t}[\sigma_{s}]^{\g}_{\gamma(s)} = \frac{\rm d}{{\rm d}s}\bigg|_{t}\big[\Fl^{\ell(X)}_{s,0}\circ\sigma\circ\Fl^{X}_{0,s}\big]_{\gamma(s)}
\end{gather*}
for all $t\in[0,\epsilon)$, which by Proposition~\ref{tangentgerm} implies the coordinate expression\vspace{-1ex}
\begin{gather*}
\frac{\rm d}{{\rm d}s}\bigg|_{t}\big(\sigma^{i}_{s}(x^{1},\dots,x^{n})- \big(\Fl^{\ell(X)}_{s,0}\circ\sigma\circ\Fl^{X}_{0,s}\big)^{i}\big(x^{1},\dots,x^{n}\big)\big) = 0.
\end{gather*}
for $i=1,\dots,m$ and for $\big(x^{1},\dots,x^{n}\big)\in\OO$. This being true for all $t\in[0,\epsilon)$, it follows that the difference $\sigma^{i}_{t}-\big(\Fl^{\ell(X)}_{t,0}\circ\sigma\circ\Fl^{X}_{0,t}\big)^{i}$ is constant in $t$ on $\OO$ for all $i$, and since $\sigma^{i}_{0} = \sigma^{i}$ on $\OO$ we obtain $[\sigma_{t}]^{\g}_{\gamma(t)} = \big[\Fl^{\ell(X)}_{t,0}\circ\sigma\circ\Fl^{X}_{0,t}\big]^{\g}_{\gamma(t)}$ for all $t\in[0,\epsilon)$. Uniqueness on $[0,d]$ now follows from the compactness of $[0,d]$.


It remains only to show smoothness of $L\big(\pi_{B}^{k,\FF}\big)$. Let $\varphi_{\PP}:=\big(\tilde{\gamma},[\tilde{X}]^{\g},\tilde{d}\big)\colon U\rightarrow\PP(\FF)$ and $\varphi_{\Gamma}:=\big(\tilde{x},[\tilde{\sigma}]^{k}_{\tilde{x}}\big)\colon V\rightarrow\Gamma_{k}(\pi_{B})^{\FF}$ be plots. We~need to show that the map from $W:=U\times_{s\circ\varphi_{\PP},\pi_{B}^{k,\FF}\circ\varphi_{\Gamma}}V\times\RB_{\geq0}$ to $\Gamma_{k}(\pi_{B})^{\FF}$ defined by
\begin{gather*}
(u,v,t)\mapsto \big(\tilde{\gamma}(u)(t),\big[\Fl^{\ell(\tilde{X}(u))}_{t,0}\circ \tilde{\sigma}(v)\circ\Fl^{\tilde{X}(u)}_{0,t}\big]^{k}_{\tilde{\gamma}(u)(t)}\big)
\end{gather*}
is smooth. The~map $\tilde{\gamma}$ is already a plot of $\PP(M)$ by definition. Recalling that $(\Gamma_{\pi_{B}})_{\loc}$ denotes the space of all locally defined sections of $\pi_{B}$ equipped with the diffeology of Proposition~\ref{functprop}, it~suffices to show that the map
\begin{gather*}
(u,v,t)\mapsto\kappa(u,v,t):=\big(\Fl^{\ell(\tilde{X}(u))}_{t,0}\circ\tilde{\sigma}(v)\circ\Fl^{\tilde{X}(u)}_{0,t}\big)\in(\Gamma_{\pi_{B}})_{\loc}
\end{gather*}
is smooth. That is, fixing $(u_{0},v_{0},t_{0})\in W$ and $x_{0}\in\dom(\kappa(u_{0},v_{0},t_{0}))$, we must find an open neighbourhood $\WW$ of $(u_{0},v_{0},t_{0})$ in $W$ and an open neighbourhood $\OO$ of $x_{0}$ in $M$ such that $\OO\subset\dom(\kappa(u,v,t))$ for all $(u,v,t)\in\WW$ and for which the map
{\samepage\begin{gather*}
\WW\times\OO\ni(u,v,t,x)\mapsto \kappa(u,v,t)(x)\in B
\end{gather*}
is smooth in the usual sense. Setting $d_{0}:=d(u_{0})$, we have the following.

}

\begin{enumerate}\itemsep=0pt
\item Using Lemma~\ref{technicallemma} together with the smoothness of $\ell$ as a map $\Gamma_{\loc}(\FF)\rightarrow(\pi_{B})_{*}(\XF_{B})_{\loc}$, we can find open neighbourhoods $\OO^{B}\ni\tilde{\sigma}(v_{0})\big(\Fl^{\tilde{X}(u_{0})}_{0,t_{0}}(x_{0})\big)$ in $B$, $\UU_{1}\ni u_{0}$ in $U$ and $I_{1}\supset[0,t_{0}]$ in $\RB_{\geq0}$ for which $\OO^{B}\subset\dom\big(\Fl^{\ell(\tilde{X}(u))}_{t,0}\big)$ for all $(u,t)\in\UU_{1}\times I_{1}$, and such that $I_{1}\times\UU_{1}\times \OO^{B}\ni(t,u,b)\mapsto\Fl^{\ell(\tilde{X}(u))}_{t,0}(b)\in B$ is smooth.
\item By definition of the diffeology on $(\Gamma_{\pi_{B}})_{\loc}$, we can find open neighbourhoods $\OO^{M}\ni\Fl^{\tilde{X}(u_{0})}_{0,t_{0}}(x_{0})$ in $M$ and $\VV\ni v_{0}$ in $V$ such that $\OO^{M}\subset\dom(\tilde{\sigma}(v))$ and $\tilde{\sigma}(v)(\OO^{M})\subset\OO^{B}$ for all $v\in\VV$, and for which $\VV\times\OO^{M}\ni(v,x)\mapsto\tilde{\sigma}(v)(x)\in B$ is smooth.
\item Again by Lemma~\ref{technicallemma}, we can find open neighbourhoods $\UU_{2}\ni u_{0}$, $I_{2}\supset[0,t_{0}]$ in $\RB_{\geq0}$ and $\OO\ni x_{0}$ in $M$ such that $\OO\subset\dom\big(\Fl^{\tilde{X}(u)}_{0,t}\big)$ and $\Fl^{\tilde{X}(u)}_{0,t}(\OO)\subset\OO^{M}$ for all $u\in\UU_{2}$ and $t\in I_{2}$, and such that $I_{2}\times\UU_{2}\times \OO\ni(t,u,x)\mapsto\Fl^{\tilde{X}(u)}_{0,t}(x)\in M$ is smooth.
\end{enumerate}
Finally, therefore, setting $\UU:=\UU_{1}\cap\UU_{2}$, $I:=I_{1}\cap I_{2}$ and
\begin{gather*}
\WW:=\UU\times_{s\circ\varphi_{\PP},\pi_{B}^{k,\FF}\circ\varphi_{\Gamma}}\VV\times I\subset W,
\end{gather*}
we have $\OO\subset\dom(\kappa(u,v,t))$ for all $(u,v,t)\in\WW$ and $\WW\times\OO\ni(u,v,t,x)\mapsto\kappa(u,v,t)(x)\in B$ is smooth.
\end{proof}
